\section*{\centering\fontsize{14}{21}\selectfont\bfseries EXECUTIVE SUMMARY}

Krisis kapasitas Tempat Pemrosesan Akhir (TPA) Bantar Gebang --- yang menerima rata-rata 7.300--7.700 ton sampah per hari dari DKI Jakarta dan telah dinyatakan \textit{overload} \parencite{dlhjakarta2026,koranjakarta2025} --- diperburuk oleh inefisiensi logistik pengangkutan sampah perkotaan yang masih bergantung pada rute dan jadwal yang statis. Armada truk sampah sering kali bergerak menjemput kontainer yang belum penuh, menyebabkan pemborosan bahan bakar, dan gagal mendistribusikan sampah secara presisi ke fasilitas pengolahan modern seperti RDF (\textit{Refuse Derived Fuel}) Plant Rorotan. \ProductName{} hadir sebagai platform \textit{Green Technology} berbasis B2B/B2G SaaS yang mengkombinasikan sensor IoT \textit{fill-level} dan \textit{Dynamic Routing AI} untuk merestrukturisasi manajemen operasional pengangkutan sampah kota.

Masalah utama pasar saat ini adalah inefisiensi logistik pengangkutan, di mana sampah diangkut berdasarkan jadwal dan bukan volume, sehingga memicu pemborosan bahan bakar armada. Penumpukan sampah yang tidak terdistribusi secara fleksibel turut mempercepat habisnya daya tampung TPA Bantar Gebang, sementara kebijakan pemilahan sampah dari sumber yang baru diberlakukan Pemprov DKI Jakarta sejak Agustus 2026 \parencite{ingub5_2026} justru membuka peluang bagi platform digital yang mampu melacak dan mengoptimalkan arus sampah terpilah tersebut.

Solusi yang kami gagas adalah sensor IoT \textit{fill-level}, yaitu sensor pintar yang dipasang pada kontainer TPS atau gedung komersial untuk memantau kapasitas sampah secara \textit{real-time}. Ketika sensor menunjukkan kapasitas suatu TPS sudah tinggi ($\geq$80\%), \textit{Dynamic Route Engine} berbasis AI akan memetakan rute pengangkutan tercepat secara otomatis. Sebelum dibuang ke TPA, \textit{Predictive RDF Allocation System} membantu mengarahkan aliran sampah bernilai kalor tinggi ke fasilitas pengolahan RDF, sehingga mengurangi beban pembuangan akhir di Bantar Gebang.

\ProductName{} mengadopsi pendekatan \textit{asset-light} (memanfaatkan armada pengangkut sampah milik pemerintah maupun swasta yang telah ada) dengan dua aliran pendapatan: \textit{Business to Government}, di mana pemerintah membayar lisensi tahunan SaaS beserta unit sensor IoT untuk pengawasan armada kota; dan \textit{Business to Business} bagi pengelola kawasan mandiri, mal, atau vendor swasta, melalui penjualan unit sensor IoT serta langganan platform bulanan untuk analisis limbah dan pelaporan keberlanjutan (ESG).

\ProductName{} bertujuan memangkas biaya operasional BBM armada pengangkut sampah, memotong waktu tempuh rute, serta meningkatkan efisiensi pengalihan sampah bernilai kalor tinggi ke fasilitas RDF. Inovasi ini mengubah tata kelola sampah analog menjadi ekosistem \textit{smart logistics} yang lebih efisien dan siap terhubung dengan agenda reformasi pengelolaan sampah DKI Jakarta.
